\documentclass[11pt,a4paper]{article}

\usepackage[T2A]{fontenc}
\usepackage[utf8]{inputenc}
\usepackage[russian]{babel}
\usepackage{lmodern}
\usepackage[margin=2cm]{geometry}
\usepackage{microtype}
\usepackage{enumitem}
\usepackage{titlesec}

\titleformat{\section}{\normalsize\bfseries}{\thesection.}{0.5em}{}
\titlespacing*{\section}{0pt}{8pt}{3pt}
\setlength{\parskip}{2pt}
\setlist[itemize]{leftmargin=1.2em, itemsep=1pt, topsep=2pt}
\pagestyle{empty}

\begin{document}

\begin{center}
  {\large\bfseries Практическая значимость проекта «Дневник калорий»}\\[3pt]
\end{center}

\section{Какую проблему решает проект}
Контроль питания нужен очень разным людям: тем, кто худеет или набирает массу,
спортсменам, людям с диетическими рекомендациями врача. Главная причина, по которой
дневник питания обычно бросают, "--- трудоёмкость: каждый продукт нужно найти в базе,
взвесить и внести вручную. Проект снижает этот барьер. Пользователь отправляет
Telegram-боту фотографию приёма пищи, система распознаёт блюда и продукты,
подставляет их калорийность и сохраняет запись в личный дневник. Вместо нескольких
минут ручного ввода требуется одно фото.

\section{Кому и чем проект полезен}
\begin{itemize}
  \item \textbf{Обычным пользователям.} Быстрый учёт питания без установки отдельного
        приложения: Telegram уже есть на телефоне.
  \item \textbf{Пользователям региона.} Словарь из 417 классов включает блюда казахской
        и центральноазиатской кухни (бешбармак, плов, самса, баурсак, кумыс и др.)
        наряду с продуктами европейского рациона из датасета MyFoodRepo. Глобальные
        приложения, как правило, покрывают такую кухню хуже.
  \item \textbf{Для анализа рациона.} Накопленная история приёмов пищи показывает
        динамику питания и может быть передана диетологу или тренеру.
  \item \textbf{Для дальнейших разработок.} Подготовленный и очищенный датасет
        (разметка в формате YOLO, единые синонимы классов) и воспроизводимый пайплайн
        можно использовать в других проектах по распознаванию еды.
\end{itemize}

\section{Ограничения}
Оценка калорийности по фотографии остаётся приближённой. Модель различает продукты,
но не знает вес порции, состав смешанных блюд, количество масла или соуса. Текущий
бейзлайн (YOLO11n, одна эпоха обучения) показывает mAP50 менее 1\% на тестовой
выборке, поэтому для реального использования пока не пригоден; эксперимент с
расширенной аугментацией на таком коротком обучении качество не улучшил. Результат
работы системы следует считать справочным и не рассматривать как медицинскую
рекомендацию.

\section{Путь к практической ценности}
Для практического применения необходимо: обучить более крупные модели на большем числе
эпох, улучшить качество на редких классах и добавить этап уточнения (пользователь
подтверждает или исправляет распознанное и указывает размер порции; по плану проекта
здесь используется LLM). Такой подход с участием человека позволяет получить
приемлемую точность учёта и превращает проект из учебного бейзлайна в работающий
инструмент, особенно востребованный там, где местная кухня плохо представлена в
существующих приложениях.

\end{document}